\begin{abstract}
本文研究分数次极大算子在尺度不变阈值 $\alpha|E|^{\beta/n}$ 下的光晕。
严格见证不等式同时约束局部密度和全局捕获比例，并给出最佳代数对称差误差。
结合近重合几何，证明非中心方体和强矩形的端点增量为 $1-\alpha$ 阶，
中心方体为 $(1-\alpha)^n$ 阶。通过共同核质量修改、面层压缩及凸纤维对称平均，
进一步确定中心方体在全体可测输入上的精确首项，并给出紧分数阶区间上的统一下一阶余项。对单位方体，联合优化中心和平移尺度，得到显式切向
多面体及其体积 $2^n/[k!(2c)^{n-k}]$，其中 $c=\min\{\beta,n-\beta\}$、
$k=\lfloor2c\rfloor$。一维情形求得高阈值精确常数；外层权情形建立集中模量的双向
控制和 $A_\infty/A_\infty^\ast$ 刻画。多线性对角模型精确约化到线性参数 $\beta/m$，
一般输入则满足几何平均归一化下的双边夹逼。最后推导全体中心限制型高矩首项、分数阶
消失时的过渡下剖面，以及 Riesz、Wolff 势的最佳方体比较和直接应用。
一维闭式还给出全阶统一余项和临界窗口的四阶展开。对任意固定满维锚定凸体，进一步建立尺度与平移联合切向极限，并求得固定方体的加权首项。方体内部的尺度优化可精确有限化；二维固定模型由此得到四分支面积公式、全部临界窗口的过渡系数，以及带显式余项界的任意阶统一边界截断表示。
任意维固定方体还具有显式次首项，以及任意阶统一径向截断；其误差阶在一般参数类中不可提高。
在中心对称半阶模型中，切向体恰为极投影体，由Petty不等式确定椭球极值。
二维半阶由精确两点差体得到显式体积上包络；独立八尺度的交并预算与各向异性侧弧传递进一步证明，该包络与真实全体常数之间的亏损阶恰为 $(1-\alpha)^3$，并给出在 $0<1-\alpha\le2^{-14}$ 时的显式减项 $2^{-68}(1-\alpha)^3$，从而使全体三阶上极限严格小于 $10/3$。矩形下界给出下极限至少 $2$；精确三阶系数及极限存在性仍待研究。进一步对与正方形对称差为 $O((1-\alpha)^2)$ 的任意输入证明统一三阶上界，四阶余项阶数不可降低；偏心参数为 $\rho$ 的有间隙近矩形类还具有三阶上界 $2+2|\rho|(1-|\rho|)\le5/2$；有界角孔具有等号刚性和定性稳定性。有限四轴见证还迫使任意输入在测度上二阶接近矩形；方向变化的共同尺度八点容量具有解析边界与显式十字形最优输入。
\end{abstract}
\noindent\textbf{关键词：}分数次极大算子；尺度不变光晕；Tauberian 常数；切向多面体；外层权；多线性传递。
